\section{Monitoring and Observability}
\label{sec:Monitoring}

\subsection{Canonical process-events topic}

Durga emits normalized lifecycle events on the logical channel \texttt{process-events}.
Generated workers, gateways, subprocess handlers, boundary handlers and helper publishers
all emit events describing process progress, waiting states, completion, failure,
escalation and cancellation. The default physical topic is
\texttt{process-events-\{processId\}} so that each process definition is isolated at the
transport level. The monitoring topology subscribes to \texttt{process-events-.*} to
consume events from all processes.

This separates execution transport from queryable state. Execution topics remain useful for
routing, but the monitoring model depends on projections built from explicit lifecycle events.
For production use, treat the lifecycle topic as an operational contract: retention,
access control, schema compatibility and replay policy determine how far back monitoring can
recover and how safely downstream tools can consume process state.

\subsection{Process self-registration}

Processes self-register with the monitoring server by publishing their BPMN 2.0 model
to the \texttt{process-models} Kafka topic. Each generated project includes a
\texttt{ModelRegistration} CDI bean that runs at startup:

\begin{itemize}
\item reads the BPMN file from the classpath (\texttt{src/main/resources/\{name\}.bpmn}),
\item publishes it to \texttt{process-models} with the process ID as the Kafka message key,
\item logs success or failure without blocking the process startup.
\end{itemize}

The monitoring server caches models in the \texttt{process-models-store} state store.
The \texttt{/api/diagram?processId=\textit{id}} endpoint serves the cached model.
Model updates overwrite previous versions (latest-wins semantics).

This registry makes processes discoverable: \texttt{/api/processes/list} merges
process IDs from both the count store (processes with live instances) and the model
store (registered processes), so newly deployed processes appear in the dashboard
before their first lifecycle event arrives.

\subsection{Kafka Streams topology}

The monitoring topology subscribes to \texttt{process-events-.*} via a regex
pattern, monitoring all processes in a single instance. State stores use
unsuffixed names. The SPA shows an inventory of all known processes with
drill-down to per-process detail.

The topology consumes lifecycle events and materializes:
\begin{itemize}
\item latest process state by instance,
\item counts by current state,
\item latency summaries,
\item stuck-instance views,
\item trend buckets,
\item BPMN model cache.
\end{itemize}

The main entry point is \texttt{org.gautelis.durga.monitoring.ProcessMonitoringApp}.

\subsection{HTTP API}

The monitoring service exposes endpoints for health, process discovery, instance
lookup, counts, latency, stuck instances, diagram serving, and a dashboard.
All monitoring paths are served under the \texttt{/api/} prefix; a root \texttt{/health} endpoint is also available.

The local development loop is:
\begin{enumerate}
\item generate a process,
\item the process auto-registers its BPMN model on boot,
\item publish scenarios or task inputs,
\item observe lifecycle events,
\item query current state and latency,
\item view the BPMN diagram with activity overlays,
\item iterate on the BPMN model or generated runtime.
\end{enumerate}

In production, the loop is stricter: monitor Streams state, query-store
readiness, consumer lag on lifecycle and monitoring topics, stuck-instance
counts, latency percentiles, DLQ volume and failed lifecycle events. The
dashboard is an operator view over the Kafka Streams read model, not a
replacement for Kafka broker, consumer group and infrastructure monitoring.

\subsection{Fault detection}

\label{sec:FaultDetection}

The fault detection topology consumes the same \texttt{process-events-.*}
stream as the monitoring topology and evaluates each lifecycle event against
a set of configured \texttt{AlarmConfig} rules. When a rule's criteria are
met, an \texttt{AlarmEvent} is published to the \texttt{fault-alarms} topic
for downstream notification, alerting, or dashboard integration.

\subsubsection{Alarm configuration}

Alarms are defined with the following scope and syndrome:

\begin{description}
\item[Scope] Each config can target a specific process (\texttt{processId}),
  a specific activity (\texttt{activityId}), both, or neither (all events).
\item[Event type] Which \texttt{ProcessEvent.EventType} triggers evaluation
  (e.g. \texttt{PROCESS\_FAILED}, \texttt{ACTIVITY\_ESCALATED}).
\item[Syndrome] How observed faults are aggregated into an alarm (see below).
\item[Severity] \texttt{WARN} or \texttt{CRITICAL}.
\item[Message] A template rendered at alarm time; supports
  \texttt{\$\{processId\}}, \texttt{\$\{activityId\}},
  \texttt{\$\{count\}}, and \texttt{\$\{processInstanceId\}}.
\end{description}

\subsubsection{Syndromes}

\begin{description}

\item[HARD\_ERROR] First occurrence triggers alarm immediately. No
threshold---every matching event produces an alarm. Suitable for critical
failures that must never happen (e.g. schema validation failure on a
compliance-critical field).

\item[COUNTED] Counts occurrences per (config, processId, activityId).
The alarm fires once when the accumulated count exceeds
\texttt{threshold}, then suppresses until the count drops back below the
threshold (re-arm). Suitable for transient errors that become significant
in volume (e.g. 5 task failures per process).

\item[SLIDING\_WINDOW] Counts occurrences within a sliding time window.
Old occurrences leak out after \texttt{windowDuration}. The alarm fires
once when the windowed count exceeds \texttt{threshold}, then suppresses
until the count drops back below threshold (re-arm). Suitable for
burst detection (e.g. more than 3 \texttt{PROCESS\_FAILED} events in a
60-second window).

\end{description}

\subsubsection{Alarm emission}

Triggered alarms are serialized as JSON and published to the
\texttt{fault-alarms} topic. Each \texttt{AlarmEvent} carries:

\begin{itemize}
\item \texttt{alarmId} --- unique per firing,
\item \texttt{configId} --- back-reference to the triggering config,
\item \texttt{syndrome}, \texttt{severity}, \texttt{message} --- from config,
\item \texttt{processId}, \texttt{processInstanceId}, \texttt{activityId} --- scope,
\item \texttt{triggerEventType}, \texttt{triggerTimestamp} --- what fired it,
\item \texttt{count} --- current accumulated count at firing,
\item \texttt{threshold} --- config threshold (0 for HARD\_ERROR).
\end{itemize}

\subsubsection{State stores}

The topology maintains three persistent state stores:

\begin{itemize}
\item \texttt{fault-counts-store} --- current count per (config, processId, activityId).
\item \texttt{fault-windows-store} --- list of epoch-millis timestamps per scope (SLIDING\_WINDOW only).
\item \texttt{fault-last-alarm-store} --- count at which the last alarm fired, used for single-fire suppression.
\end{itemize}

\subsubsection{Example configuration}

\begin{verbatim}
AlarmConfig hard = new AlarmConfig("schema-fail-critical", null,
    "validate_high_value", ProcessEvent.EventType.ACTIVITY_ESCALATED,
    AlarmSyndrome.HARD_ERROR, 0, null,
    AlarmSeverity.CRITICAL, "Schema validation escalated");

AlarmConfig counted = new AlarmConfig("proc-fail-warn", "order-proc",
    null, ProcessEvent.EventType.PROCESS_FAILED,
    AlarmSyndrome.COUNTED, 5, null,
    AlarmSeverity.WARN, "${count} failures in ${processId}");

AlarmConfig window = new AlarmConfig("burst-detect", null, null,
    ProcessEvent.EventType.PROCESS_FAILED,
    AlarmSyndrome.SLIDING_WINDOW, 3, Duration.ofMinutes(1),
    AlarmSeverity.CRITICAL, "Burst: ${count} failures in 1 min");
\end{verbatim}

\subsubsection{Integration with the monitoring server}

The fault detection topology runs alongside the monitoring topology in the
same Kafka Streams application. Both share the \texttt{process-events-.*}
subscription. The \texttt{FaultDetectionTopology.buildTopology(\dots)}
method returns a standalone \texttt{Topology} that can be started in the
same \texttt{KafkaStreams} instance or in a separate one with its own
application ID.

Unit tests are in \texttt{FaultDetectionTest} (13 scenarios); see
Appendix~\ref{sec:IntegrationTestAppendix} for the test catalogue.
